\appendixheading

*引理 1. --- 令 X、B 为两个局部紧空间，\(\pi\) 为从 X 到 B 的映射，\(\nu\) 为 B 上的正测度。令 \(b \mapsto \lambda_b\)（\(b \in \mathbf{B}\)）为 X 上正测度的一个 \(\nu\)-适当族，使得对每个 \(b \in \mathbf{B}\)，测度 \(\lambda_b\) 集中在 \(\overline{\pi}^1(b)\) 上。置 \(\mu = \int \lambda_b d\nu(b)\)，并假设映射 \(\pi\) 是 \(\mu\)-可测的。

\begin{enumerate}
\def\labelenumi{\alph{enumi})}
\item
  若 \(\mathbf{N} \subset \mathbf{B}\) 是局部 \(\nu\)-零测的，则 \(\overline{\pi}^1(\mathbf{N})\) 是局部 \(\mu\)-零测的。
\item
  若 \(f\) 是 \(\nu\)-可测函数（定义在 \(\mathbf{B}\) 上，值域为某个拓扑空间），则 \(f \circ \pi\) 是 \(\mu\)-可测函数（定义在 \(\mathbf{X}\) 上）。
\end{enumerate}

令 K 为 X 的一个紧子集。我们要证明 \(\overline{\pi}^{1}(N) \cap K\) 是 \(\mu\)-零测的，并且 \(f \circ \pi\) 在 K 上的限制是 \(\mu\)-可测的。现在，K 是一个 \(\mu\)-零测集与一列紧子集 \(K_{n}\) 的并，且 \(\pi|K_{n}\) 连续。只需证明 \(\overline{\pi}^{1}(N) \cap K_{n}\) 是 \(\mu\)-零测的，并且 \(f \circ \pi\) 在 \(K_{n}\) 上的限制是 \(\mu\)-可测的。因此以下不妨假设 \(\pi|K\) 连续。于是 \(\pi(K) = K'\) 是紧的。由于 \(\overline{\pi}^{1}(N) \cap K = \overline{\pi}^{1}(N \cap K') \cap K\)，且 \(N \cap K'\) 是 \(\nu\)-零测的，以下不妨假设 N 是 \(\nu\)-零测的。于是 N 包含于一个 \(\nu\)-零测集 \(N'\) 中，而它是开集的可数交（Ch. IV, §4, No.~6, Th. 4 的 Cor. 2）。由于 \(\pi\) 是 \(\mu\)-可测的，\(\overline{\pi}^{1}(N')\) 是 \(\mu\)-可测的（Ch. IV, §5, No.~5, Prop. 7），所以 \(\overline{\pi}^{1}(N') \cap K\) 是 \(\mu\)-可积的，并且（Ch. V, §3, No.~4, Th. 1）

\[
\mu \big (\overline {{{{\pi}}}} ^ {1} (\mathbf {N} ^ {\prime}) \cap \mathbf {K} \big) = \int_ {\mathbf {B}} \lambda_ {b} \big (\overline {{{{\pi}}}} ^ {1} (\mathbf {N} ^ {\prime}) \cap \mathbf {K} \big) d \nu (b).
\]

现在，若 \(b \notin \mathbf{N}'\)，则集合 \(\overline{\pi}(\mathbf{N}') \cap \mathbf{K}\) 对 \(\lambda_b\) 是零测的，因为 \(\lambda_b\) 集中在 \(\overline{\pi}(b)\) 上。因此 \(\mu\big(\overline{\pi}(\mathbf{N}') \cap \mathbf{K}\big) = 0\)。从而，\(\overline{\pi}(\mathbf{N}) \cap \mathbf{K}\) 是 \(\mu\)-零测的，a) 已得证。另一方面，存在 \(\mathbf{K}'\) 的一个划分，由一个 \(\nu\)-零测集 \(\mathbf{M}\) 和一列紧集 \((\mathbf{K}_n')\) 构成，使得 \(f|_{\mathbf{K}_n'}\) 连续。于是 \(f \circ \pi\) 在每个集合 \(\overline{\pi}^1(\mathbf{K}) \cap \mathbf{K}_n'\) 上的限制是连续的，而 \(\overline{\pi}^1(\mathbf{M}) \cap \mathbf{K}\) 根据 a) 是 \(\mu\)-零测的，因此 \(f \circ \pi\) 在 \(\mathbf{K}\) 上的限制确实是 \(\mu\)-可测的。

